% GENERATED FILE. Edit canonical lessons, then run npm run generate:books.
% canonical-source-hash: fnv1a64:c50179a0715bc418
% canonical-lessons: ML-C389-grhapravesam, ML-C389-namakaranam, ML-C389-madhuvidhu, ML-C389-tamasamakkuka, ML-C389-kutiyeruka

\chapter{New Homes and New Names}
\label{ch:ml-new-homes-and-new-names}
\hlchaptermodality{389}

\begin{chapteropening}
\textbf{By the end of this chapter:} \emph{I can say a housewarming, a naming ceremony, a honeymoon, to settle and to migrate.}

Complete the last lesson of chapter 389: I can say a housewarming, a naming ceremony, a honeymoon, to settle and to migrate.
\end{chapteropening}

\section[\texorpdfstring{\ml{ഗൃഹപ്രവേശം} (gṛhapravēśaṁ)}{ഗൃഹപ്രവേശം (gṛhapravēśaṁ)}]{\ml{ഗൃഹപ്രവേശം} (gṛhapravēśaṁ) — a housewarming}

\label{lesson:ML-C389-grhapravesam}



\begin{quote}
Before the new one: say the Malayalam for to congratulate, then the Malayalam for to bestow.
\end{quote}

\subsection*{You'll want to know: \ml{ഗൃഹപ്രവേശം}}
\textbf{\ml{ഗൃഹപ്രവേശം}} — \emph{gṛhapravēśaṁ} — \textquotedblleft{}a housewarming\textquotedblright{}.

\begin{quote}
\ml{അവരുടെ} \ml{പുതിയ} \ml{വീടിന്റെ} \ml{ഗൃഹപ്രവേശം} \ml{ഞായറാഴ്ച} \ml{ആണ്}.
\end{quote}

\begin{quote}
\emph{avaruṭe putiya vīṭinṟe gṛhapravēśaṁ ñāyaṟāḻca āṇŭ.}
\end{quote}

\begin{quote}
Their new house's housewarming is on Sunday.
\end{quote}

Literally entering the house: \ml{ഗൃഹം}, a house, + \ml{പ്രവേശം}, entry.

\begin{cousinweb}
\textbf{\ml{ഗൃഹപ്രവേശം}} (\emph{gṛhapravēśaṁ}) is borrowed from Sanskrit \emph{gṛhapraveśa}, \textquotedblleft{}solemn entrance into a house\textquotedblright{} (Monier-Williams), from \emph{gṛha}, \textquotedblleft{}a house\textquotedblright{}, and \emph{praveśa}, \textquotedblleft{}entering\textquotedblright{}.
\end{cousinweb}

\begin{grammarlens}[title={What you've built}]
One more word for telling what happened.
\end{grammarlens}

\subsection*{Your turn}
\begin{itemize}
\raggedright
  \item \emph{Say it:} \emph{gṛhapravēśaṁ}
  \item \emph{Say it:} \emph{gṛhapravēśaṁ}, once more
  \item \emph{Say it:} \emph{avaruṭe putiya vīṭinṟe gṛhapravēśaṁ ñāyaṟāḻca āṇŭ.} — Their new house's housewarming is on Sunday.
  \item \emph{Recall:} say the Malayalam for to congratulate, then the Malayalam for to bestow, then say \emph{gṛhapravēśaṁ} again
\end{itemize}

\begin{tcolorbox}[breakable,colback=teal!4,colframe=teal!35!black,title={Before you move on}]
What does \ml{ഗൃഹപ്രവേശം} mean? (\textquotedblleft{}A housewarming\textquotedblright{}.) Say it once more.
\end{tcolorbox}
\section[\texorpdfstring{\ml{നാമകരണം} (nāmakaraṇaṁ)}{നാമകരണം (nāmakaraṇaṁ)}]{\ml{നാമകരണം} (nāmakaraṇaṁ) — a naming ceremony}

\label{lesson:ML-C389-namakaranam}



\begin{quote}
Before the new one: say the Malayalam for to bestow, then the Malayalam for a housewarming.
\end{quote}

\subsection*{You'll want to know: \ml{നാമകരണം}}
\textbf{\ml{നാമകരണം}} — \emph{nāmakaraṇaṁ} — \textquotedblleft{}a naming ceremony\textquotedblright{}.

\begin{quote}
\ml{കുഞ്ഞിന്റെ} \ml{നാമകരണം} \ml{ഇരുപത്തിയെട്ടാം} \ml{ദിവസം} \ml{ആയിരുന്നു}.
\end{quote}

\begin{quote}
\emph{kuññinṟe nāmakaraṇaṁ irupattiyeṭṭāṁ divasaṁ āyirunnu.}
\end{quote}

\begin{quote}
The baby's naming was on the twenty-eighth day.
\end{quote}

Also the naming of anything new, a road or a building.

\begin{cousinweb}
\textbf{\ml{നാമകരണം}} (\emph{nāmakaraṇaṁ}) is borrowed from Sanskrit \emph{nāmakaraṇa}, \textquotedblleft{}the ceremony of naming after birth\textquotedblright{} (CDIAL 7066), from \emph{nāman}, \textquotedblleft{}a name\textquotedblright{}, and \emph{karaṇa}, \textquotedblleft{}a making\textquotedblright{}.
\end{cousinweb}

\begin{grammarlens}[title={What you've built}]
One more word for telling what happened.
\end{grammarlens}

\subsection*{Your turn}
\begin{itemize}
\raggedright
  \item \emph{Say it:} \emph{nāmakaraṇaṁ}
  \item \emph{Say it:} \emph{nāmakaraṇaṁ}, once more
  \item \emph{Say it:} \emph{kuññinṟe nāmakaraṇaṁ irupattiyeṭṭāṁ divasaṁ āyirunnu.} — The baby's naming was on the twenty-eighth day.
  \item \emph{Recall:} say the Malayalam for to bestow, then the Malayalam for a housewarming, then say \emph{nāmakaraṇaṁ} again
\end{itemize}

\begin{tcolorbox}[breakable,colback=teal!4,colframe=teal!35!black,title={Before you move on}]
What does \ml{നാമകരണം} mean? (\textquotedblleft{}A naming ceremony\textquotedblright{}.) Say it once more.
\end{tcolorbox}
\section[\texorpdfstring{\ml{മധുവിധു} (madhuvidhu)}{മധുവിധു (madhuvidhu)}]{\ml{മധുവിധു} (madhuvidhu) — a honeymoon}

\label{lesson:ML-C389-madhuvidhu}



\begin{quote}
Before the new one: say the Malayalam for a housewarming, then the Malayalam for a naming ceremony.
\end{quote}

\subsection*{You'll want to know: \ml{മധുവിധു}}
\textbf{\ml{മധുവിധു}} — \emph{madhuvidhu} — \textquotedblleft{}a honeymoon\textquotedblright{}.

\begin{quote}
\ml{വിവാഹം} \ml{കഴിഞ്ഞ്} \ml{അവർ} \ml{മധുവിധുവിന്} \ml{ഒരു} \ml{ദ്വീപിൽ} \ml{പോയി}.
\end{quote}

\begin{quote}
\emph{vivāhaṁ kaḻiññŭ avar madhuvidhuvinŭ oru dvīpil pōyi.}
\end{quote}

\begin{quote}
After the wedding they went to an island for their honeymoon.
\end{quote}

Usually explained as a calque on the English word: \ml{മധു}, honey, + \ml{വിധു}, the moon.

\begin{cousinweb}
\textbf{\ml{മധുവിധു}} (\emph{madhuvidhu}) is a modern compound built from two Sanskrit words, \emph{madhu}, \textquotedblleft{}honey\textquotedblright{} (CDIAL 9784), and \emph{vidhu}, a name of the moon (Monier-Williams). It is usually explained as a calque on English \emph{honeymoon}.
\end{cousinweb}

\begin{grammarlens}[title={What you've built}]
One more word for telling what happened.
\end{grammarlens}

\subsection*{Your turn}
\begin{itemize}
\raggedright
  \item \emph{Say it:} \emph{madhuvidhu}
  \item \emph{Say it:} \emph{madhuvidhu}, once more
  \item \emph{Say it:} \emph{vivāhaṁ kaḻiññŭ avar madhuvidhuvinŭ oru dvīpil pōyi.} — After the wedding they went to an island for their honeymoon.
  \item \emph{Recall:} say the Malayalam for a housewarming, then the Malayalam for a naming ceremony, then say \emph{madhuvidhu} again
\end{itemize}

\begin{tcolorbox}[breakable,colback=teal!4,colframe=teal!35!black,title={Before you move on}]
What does \ml{മധുവിധു} mean? (\textquotedblleft{}A honeymoon\textquotedblright{}.) Say it once more.
\end{tcolorbox}
\section[\texorpdfstring{\ml{താമസമാക്കുക} (tāmasamākkuka)}{താമസമാക്കുക (tāmasamākkuka)}]{\ml{താമസമാക്കുക} (tāmasamākkuka) — to settle; to take up residence}

\label{lesson:ML-C389-tamasamakkuka}



\begin{quote}
Before the new one: say the Malayalam for a naming ceremony, then the Malayalam for a honeymoon.
\end{quote}

\subsection*{You'll want to know: \ml{താമസമാക്കുക}}
\textbf{\ml{താമസമാക്കുക}} — \emph{tāmasamākkuka} — \textquotedblleft{}to settle; to take up residence\textquotedblright{}.

\begin{quote}
\ml{അവർ} \ml{ഇപ്പോൾ} \ml{നഗരത്തിൽ} \ml{താമസമാക്കി}.
\end{quote}

\begin{quote}
\emph{avar ippōḷ nagarattil tāmasamākki.}
\end{quote}

\begin{quote}
They have now settled in the city.
\end{quote}

Past: \ml{താമസമാക്കി}, \emph{tāmasamākki}.

\ml{താമസം}, a stay, + \ml{ആക്കുക}, to make: to make a place your home.

\begin{cousinweb}
\textbf{\ml{താമസമാക്കുക}} (\emph{tāmasamākkuka}) joins \ml{താമസം}, \textquotedblleft{}a stay, a residence\textquotedblright{}, usually traced to Sanskrit \emph{tāmasa}, \textquotedblleft{}dark; dull, slow\textquotedblright{} (CDIAL 5775), as is \ml{താമസിക്കുക} (\emph{tāmasikkuka}), \textquotedblleft{}to stay\textquotedblright{}, and \ml{ആക്കുക}, \textquotedblleft{}to make\textquotedblright{}.
\end{cousinweb}

\begin{grammarlens}[title={What you've built}]
One more word for telling what happened.
\end{grammarlens}

\subsection*{Your turn}
\begin{itemize}
\raggedright
  \item \emph{Say it:} \emph{tāmasamākkuka}
  \item \emph{Say it:} \emph{tāmasamākkuka}, once more
  \item \emph{Say it:} \emph{avar ippōḷ nagarattil tāmasamākki.} — They have now settled in the city.
  \item \emph{Recall:} say the Malayalam for a naming ceremony, then the Malayalam for a honeymoon, then say \emph{tāmasamākkuka} again
\end{itemize}

\begin{tcolorbox}[breakable,colback=teal!4,colframe=teal!35!black,title={Before you move on}]
What does \ml{താമസമാക്കുക} mean? (\textquotedblleft{}To settle; to take up residence\textquotedblright{}.) Say it once more.
\end{tcolorbox}
\section[\texorpdfstring{\ml{കുടിയേറുക} (kuṭiyēṟuka)}{കുടിയേറുക (kuṭiyēṟuka)}]{\ml{കുടിയേറുക} (kuṭiyēṟuka) — to migrate; to settle in a new land}

\label{lesson:ML-C389-kutiyeruka}



\begin{quote}
Before the new one: say the Malayalam for a honeymoon, then the Malayalam for to settle.
\end{quote}

\subsection*{You'll want to know: \ml{കുടിയേറുക}}
\textbf{\ml{കുടിയേറുക}} — \emph{kuṭiyēṟuka} — \textquotedblleft{}to migrate; to settle in a new land\textquotedblright{}.

\begin{quote}
\ml{പണ്ട്} \ml{ഒരുപാട്} \ml{കുടുംബങ്ങൾ} \ml{വിദേശത്തേക്ക്} \ml{കുടിയേറി}.
\end{quote}

\begin{quote}
\emph{paṇṭŭ orupāṭŭ kuṭumbaṅṅaḷ vidēśattēkkŭ kuṭiyēṟi.}
\end{quote}

\begin{quote}
Long ago many families migrated abroad.
\end{quote}

Past: \ml{കുടിയേറി}, \emph{kuṭiyēṟi}.

\ml{കുടി}, a home, + \ml{ഏറുക}, to climb in: to set up home in a new place.

\begin{cousinweb}
\textbf{\ml{കുടിയേറുക}} (\emph{kuṭiyēṟuka}) is a native Dravidian compound of \ml{കുടി}, \textquotedblleft{}a house, a hut, a family\textquotedblright{} (DEDR 1655), and \ml{ഏറുക}, \textquotedblleft{}to climb, to go up into\textquotedblright{}.
\end{cousinweb}

\begin{grammarlens}[title={What you've built}]
One more word for telling what happened.
\end{grammarlens}

\subsection*{Your turn}
\begin{itemize}
\raggedright
  \item \emph{Say it:} \emph{kuṭiyēṟuka}
  \item \emph{Say it:} \emph{kuṭiyēṟuka}, once more
  \item \emph{Say it:} \emph{paṇṭŭ orupāṭŭ kuṭumbaṅṅaḷ vidēśattēkkŭ kuṭiyēṟi.} — Long ago many families migrated abroad.
  \item \emph{Recall:} say the Malayalam for a honeymoon, then the Malayalam for to settle, then say \emph{kuṭiyēṟuka} again
\end{itemize}

\begin{tcolorbox}[breakable,colback=teal!4,colframe=teal!35!black,title={Before you move on}]
What does \ml{കുടിയേറുക} mean? (\textquotedblleft{}To migrate; to settle in a new land\textquotedblright{}.) Say it once more.
\end{tcolorbox}
